\section{Statistical Considerations}

The statistical framework for this study is that of a small, intensively
monitored, single-arm, first-in-human investigation in which the primary purpose
is to characterize early safety and to identify a recommended dose, not to test a
comparative efficacy hypothesis. Accordingly, the analyses are predominantly
descriptive and estimation-based rather than confirmatory, and the design is
governed by its operating characteristics rather than by formal statistical power.
All analyses follow the principles of the International Council for Harmonisation
(ICH) E9 guideline on statistical principles for clinical trials, are prespecified
in this section and in the Statistical Analysis Plan (SAP) that this section
summarizes, and are reported in accordance with the TRIPOD+AI statement for the
model-assisted advisory component \cite{Collins2024tripodai}.

\subsection{Statistical Hypotheses}

The study has one primary safety hypothesis and one dose-finding objective, each
matched to an analytic strategy appropriate for the small sample.

The primary safety hypothesis is framed as a one-sided assessment against a
prespecified unacceptable threshold. Let $\pi$ denote the true incidence of
device-related or procedure-related serious adverse events (serious AEs) within
the 30-day perioperative window. The null and alternative hypotheses are
$H_0\!: \pi \geq \pi_0$ versus $H_1\!: \pi < \pi_0$, where $\pi_0$ is the
prespecified unacceptable serious-AE rate. The trial is designed so that a
device-related serious-AE rate at or above the cohort-level halt threshold defined
in \S9.4.6 (a rate $\geq$ one third within a cohort, or two device-related deaths)
triggers a mandatory data and safety monitoring board (DSMB) review and a possible
halt. Because of the small sample, this hypothesis is evaluated by estimation: the
observed serious-AE proportion is reported with an exact (Clopper-Pearson) two-sided
$95\%$ confidence interval, and acceptability is judged by whether the interval
excludes, or fails to exclude, the unacceptable threshold $\pi_0$, rather than by a
fixed-rejection significance test that a single-arm safety cohort cannot reliably
support.

The dose-finding objective is to identify the maximum tolerated dose (MTD) and the
recommended Phase 2 dose (RP2D) of perioperative daraxonrasib using the standard
3+3 rule (\S4.1, \S6). This objective is operational rather than hypothesis-testing:
the MTD is the highest dose level at which fewer than two of up to six participants
experience a dose-limiting toxicity (DLT), and the RP2D is selected by integrating
the MTD with the totality of the safety, tolerability, telemetry, and available
pharmacodynamic information. No formal type I error is spent on dose selection;
the operating characteristics of the 3+3 rule (\S9.2) govern the dose-finding
behavior. ICH E9 estimands principles are applied to the supportive efficacy
endpoints: the population of interest, the endpoint, the handling of
intercurrent events (for example open conversion, withdrawal, or death), and the
population-level summary are each specified for every endpoint in \S9.4.

\subsection{Sample Size Determination}

The planned sample size is up to $n=18$ treated participants, allocated across
three daraxonrasib dose levels (DL1 through DL3) by the 3+3 rule. This number is
not derived from a power calculation, because the study tests no confirmatory
efficacy hypothesis; it is determined jointly by the dose-finding design and by
the early-feasibility precedent for significant-risk devices, and it is justified
by its operating characteristics.

The 3+3 rule enrolls cohorts of three. A dose level that yields no DLT in three
participants escalates; a dose level that yields one DLT in three expands to six,
and a dose level that yields two or more DLTs in up to six is exceeded, defining
the dose below it as the MTD. The maximum exposure at any single dose level is
therefore six participants, and three escalation levels give a design maximum of
$3 \times 6 = 18$ treated participants; many realizations terminate with fewer.
This cohort size is consistent with the FDA early-feasibility guidance for
significant-risk devices, under which limited initial clinical experience,
generally fewer than ten participants per device-feasibility cohort, is sought
when nonclinical simulation alone cannot supply the information needed to advance
development \cite{fda2013earlyfeasibility}. The combined ceiling of 18 thus
respects the early-feasibility convention at the device-cohort level while
providing the per-dose exposure that the drug-arm 3+3 escalation requires.

Rather than power, the design is characterized by its operating characteristics.
The probability that the 3+3 rule escalates past a given dose level is a function
of the true DLT probability at that level: a dose with a true DLT rate of roughly
$10\%$ is very likely to be tolerated and escalated, whereas a dose with a true
DLT rate of $40\%$ or higher is very likely to be identified as excessively toxic
and to define the MTD at the level below. Precision of estimation at the planned
sample is summarized in Table~\ref{tab:opchar}: with $n$ between 3 and 18, an
observed proportion is reported with an exact $95\%$ confidence interval whose
width reflects the small sample, so that, for example, zero events in 18
participants yields an exact upper $95\%$ confidence bound near $18\%$, and one
event in 18 yields an interval excluding rates above roughly one quarter. No
formal significance level is set for the dose-finding analysis; for the supportive
estimation analyses, two-sided $95\%$ confidence intervals are reported throughout
(nominal $\alpha = 0.05$, two-sided), and no multiplicity adjustment is applied
because these analyses are descriptive and non-confirmatory.

\begin{table}[htbp]
\centering
\footnotesize
\caption{Operating characteristics and exact-interval precision at the planned
sample sizes. Confidence intervals are two-sided exact (Clopper-Pearson) bounds
for a single proportion; the dose-finding behavior is that of the standard 3+3
rule. No statistical power is computed because the study tests no confirmatory
efficacy hypothesis.}
\label{tab:opchar}
\begin{tabularx}{\textwidth}{L{2.0cm} L{2.6cm} Y}
\toprule
\textbf{Quantity} & \textbf{Planned value} & \textbf{Operating characteristic / precision} \\
\midrule
Cohort size & 3, expand to 6 & 3+3 rule: escalate on 0/3; expand on 1/3; exceed dose on $\geq 2/6$. \\
Per-dose maximum & 6 & Maximum exposure at any one dose level before the MTD is declared. \\
Dose levels & 3 (DL1 to DL3) & Three-level escalation; design maximum $3 \times 6 = 18$ treated. \\
Total treated & up to 18 & At or below the early-feasibility convention of generally fewer than 10 per device cohort, summed across dose levels \cite{fda2013earlyfeasibility}. \\
0 events / 3 & point estimate $0\%$ & Exact two-sided $95\%$ CI approximately $0\%$ to $71\%$. \\
0 events / 6 & point estimate $0\%$ & Exact two-sided $95\%$ CI approximately $0\%$ to $46\%$. \\
0 events / 18 & point estimate $0\%$ & Exact two-sided $95\%$ CI approximately $0\%$ to $18\%$. \\
1 event / 18 & point estimate $5.6\%$ & Exact two-sided $95\%$ CI approximately $0.1\%$ to $27\%$. \\
Significance level & $\alpha = 0.05$ (two-sided) & Applies to supportive estimation only; no $\alpha$ is spent on dose selection. \\
\bottomrule
\end{tabularx}
\end{table}

\subsection{Populations for Analyses}

Four analysis populations are defined and applied consistently across all
analyses. The Safety population comprises all participants who received any study
intervention (any dose of daraxonrasib or any patient-facing robotic activity) and
is the primary population for all safety and primary-endpoint analyses. The
DLT-evaluable population comprises all dosed participants who either completed the
28-day DLT-evaluation window or experienced a DLT within it, and is the population
on which the 3+3 escalation decisions and the MTD determination are based. The
Per-Protocol population comprises participants who received the robotic Whipple
and the LLM-directed advisory per protocol without a major protocol deviation that
could affect the efficacy assessment, and supports the supportive surgical-quality
analyses. The modified intention-to-treat (mITT) population comprises all dosed
participants who contributed at least one post-baseline assessment, and is used
for the supportive and exploratory efficacy summaries. Figure~22 shows the
derivation of the four populations and the cohort-level interim stopping logic.

\begin{mermaidfig}
\node[mmin]  (enr) at (0,0)      {Enrolled and\\dose-assigned\\n = 18};
\node[mmstep](saf) at (4.2,1.3)  {Safety population\\all who received\\any intervention};
\node[mmstep](dlt) at (4.2,-0.0) {DLT-evaluable\\completed 28-day\\window or had DLT};
\node[mmstep](pp)  at (4.2,-1.3) {Per-protocol\\Whipple + advisory\\per protocol};
\node[mmstep](mitt)at (4.2,-2.6) {modified ITT\\dosed with $\geq 1$\\post-baseline assessment};
\node[mmdec] (stop)at (8.6,1.3)  {Interim safety\\review after each\\cohort + DSMB};
\node[mmgoal](halt)at (8.6,-0.6) {Halt: $\geq 33\%$\\device-related\\serious AE, or 2\\device-related deaths};
\node[mmgoal](cont)at (8.6,-2.6) {Continue\\escalation};
\draw[mmedge] (enr) -- (saf);
\draw[mmedge] (enr) -- (dlt);
\draw[mmedge] (enr) -- (pp);
\draw[mmedge] (enr) -- (mitt);
\draw[mmedge] (saf) -- (stop);
\draw[mmedge] (stop) -- node[right,font=\scriptsize]{met} (halt);
\draw[mmedge] (stop.south) -- node[right,font=\scriptsize,pos=0.6]{below} (cont.north);
\end{mermaidfig}
\figcaption{Figure 22. The four analysis populations (safety, dose-limiting-toxicity
evaluable, per-protocol, modified intention-to-treat) and the cohort-level interim
safety review with prespecified DSMB halt rules (a device-related serious-AE rate
at or above one third within a cohort, or two device-related deaths).}

\subsection{Statistical Analyses}

\subsubsection{General Approach}

All analyses are descriptive and estimation-based. Continuous variables are
summarized by the number of non-missing observations, mean, standard deviation,
median, minimum, and maximum; categorical variables are summarized by counts and
percentages. Because the sample is small, every proportion (incidence, rate, or
fraction) is accompanied by an exact (Clopper-Pearson) two-sided $95\%$ confidence
interval rather than a normal-approximation (Wald) interval, which is unreliable
near the bounds and at small $n$. Time-to-event variables are summarized by the
Kaplan-Meier method with the number at risk, and no inferential model is fitted
unless prespecified as exploratory. Missing data are not imputed for the primary
or safety analyses; the analysis population denominators are reported explicitly,
and missingness is described. No interim or final hypothesis test is conducted for
efficacy, so no multiplicity adjustment is required; the few inferential summaries
that appear are flagged as exploratory and non-confirmatory. All analyses are
performed in a validated, version-controlled statistical environment under a
deterministic random seed, and the analysis code and derived data sets are
deposited so that every reported figure is reproducible.

\subsubsection{Analysis of the Primary Endpoint(s)}

The co-primary safety and feasibility endpoints are analyzed on the Safety
population. The primary safety endpoint, the incidence of device-related or
procedure-related serious AEs within 30 days, is reported as a proportion with an
exact two-sided $95\%$ confidence interval, overall and by dose level, and the
upper confidence bound is compared against the prespecified unacceptable threshold
$\pi_0$ (\S9.1). The dose-finding endpoint, the MTD and the RP2D, is determined on
the DLT-evaluable population by applying the 3+3 rule to the observed DLT counts
at each dose level; the MTD is reported as the highest dose with fewer than two
DLTs in up to six participants, and the RP2D is reported with the rationale that
integrated the MTD with the safety, tolerability, and telemetry data. The
operative feasibility endpoint, the proportion of cases in which the assigned
operative tasks are completed under the staged-autonomy model without conversion
attributable to system failure, is reported as a proportion with an exact $95\%$
confidence interval. No formal significance test is applied to any primary
endpoint; acceptability is judged by the confidence intervals and the prespecified
thresholds.

\subsubsection{Analysis of the Secondary Endpoint(s)}

Secondary surgical-quality and short-term oncologic endpoints are analyzed on the
Per-Protocol population (with the Safety population as a sensitivity analysis) and
each is reported as a proportion with an exact two-sided $95\%$ confidence
interval, accompanied by the contemporaneous Dutch 2025 nationwide
robotic-pancreatoduodenectomy benchmark as a descriptive, non-inferential
comparator \cite{DutchCohort2025}. The prespecified secondary endpoints and their
comparators are summarized in Table~\ref{tab:secondary}. The margin-negative (R0)
resection rate, the rate of clinically relevant (ISGPS grade B or C) postoperative
pancreatic fistula graded per the 2016 ISGPS update \cite{Bassi2017}, the rate of
Clavien-Dindo grade III or higher complications, and 30-day and 90-day mortality
are each estimated with exact intervals. Because the comparison with the Dutch
benchmark is descriptive and the study is single-arm, no hypothesis test of
non-inferiority or superiority is performed; the benchmark is presented as context
for interpreting the observed point estimates and intervals.

\begin{table}[htbp]
\centering
\footnotesize
\caption{Secondary surgical-quality and short-term oncologic endpoints, the
analysis population, and the descriptive Dutch 2025 nationwide
robotic-pancreatoduodenectomy comparator. Comparisons are descriptive only; the
single-arm design supports estimation, not inferential between-group testing.}
\label{tab:secondary}
\begin{tabularx}{\textwidth}{L{4.0cm} L{2.8cm} Y}
\toprule
\textbf{Secondary endpoint} & \textbf{Population / method} & \textbf{Reporting and comparator} \\
\midrule
Margin-negative (R0) resection rate & Per-protocol; exact $95\%$ CI & Proportion with exact CI; descriptive comparison with the Dutch 2025 cohort \cite{DutchCohort2025}. \\
Clinically relevant pancreatic fistula (ISGPS grade B/C) & Per-protocol; exact $95\%$ CI & Graded per the 2016 ISGPS update \cite{Bassi2017}; benchmarked descriptively to Dutch 2025. \\
Clavien-Dindo grade III+ complications & Safety; exact $95\%$ CI & Proportion with exact CI; Dutch 2025 descriptive comparator. \\
30-day mortality & Safety; exact $95\%$ CI & Proportion with exact CI; Dutch 2025 descriptive comparator. \\
90-day mortality & Safety; exact $95\%$ CI & Proportion with exact CI; Dutch 2025 descriptive comparator. \\
Estimated blood loss, operative time, length of stay & Per-protocol; descriptive & Mean / median with range; Dutch 2025 phase-of-learning-curve context. \\
\bottomrule
\end{tabularx}
\end{table}

\subsubsection{Safety Analyses}

Safety is the focus of the study, and all safety analyses are performed on the
Safety population. Adverse events are coded to the Medical Dictionary for
Regulatory Activities (MedDRA) and tabulated by system organ class and preferred
term, by maximum severity using the Common Terminology Criteria for Adverse
Events (CTCAE), and by relationship to the drug, to the device, and to the
procedure. Serious AEs, DLTs, and deaths are listed individually. Treatment-emergent
AEs are summarized overall and by dose level, with the number of participants
reporting each event, the number of events, and the corresponding proportions with
exact $95\%$ confidence intervals.

Device telemetry constitutes a second, parallel safety stream unique to this
trial class. The cyber-physical record is summarized per procedure and in
aggregate: emergency-stop activations and the measured stop latency against the
$\leq 500$ ms requirement; positional and force-tolerance excursions against the
$\leq 2$ mm and $\leq 0.5$ N envelopes; vascular-exclusion gating events; and the
counts of LLM advisories that were accepted, blocked, or escalated to the human
operator by the verification-before-generation gate surface. A dedicated Physical
AI adverse-event tabulation records, for every event, whether an LLM advisory, a
robot command, a safety-limit trip, or a human override was implicated, so that
events arising from the autonomy-bearing layer are visible and traceable to a
deterministic seed and a repository commit through the hash-chained audit trail
(\S10.3). Laboratory data, vital signs, and physical-examination findings are
summarized descriptively with shift tables relative to baseline.

\subsubsection{Baseline Descriptive Statistics}

Baseline characteristics are summarized descriptively for each analysis
population. Reported variables include demographics (age, sex, race, ethnicity);
disease staging and resectability category (resectable versus borderline-resectable);
the specific KRAS G12 mutation; baseline CA 19-9; Eastern Cooperative Oncology
Group (ECOG) performance status; prior neoadjuvant therapy and number of cycles;
and relevant comorbidity and organ-function indices. No formal between-group
comparison of baseline characteristics is performed, consistent with the
single-arm design; baseline balance across dose levels is described to support
interpretation of the safety and dose-finding results.

\subsubsection{Planned Interim Analyses}

Formal interim safety analyses are conducted at the cohort level. After each
participant completes the 28-day DLT-evaluation window and before the next
participant in the cohort undergoes surgery, the accumulating safety data are
reviewed under the staggered (sentinel) enrollment rule (\S4.1); and after each
dose-level cohort is complete, the DSMB convenes for a formal interim safety
review and an escalation, expansion, hold, or halt decision. The prespecified
halt rules are quantitative and binding: enrollment and escalation halt for a
mandatory DSMB review if the device-related or procedure-related serious-AE rate
reaches or exceeds one third within a cohort, or if two device-related deaths
occur at any point. Additional pause triggers include any unexpected
device-related serious AE, any breach of a hard cyber-physical safety limit
(emergency-stop latency, positional, or force envelope), and any failure of the
pre-procedure safety-matrix verification. The DSMB may halt, pause, modify, or
permit continuation; a halt for an immediate hazard is implemented at once under
\S312.30(h) and \S10.10. No interim efficacy analysis and no efficacy-based
stopping rule are specified, so no efficacy-related type I error is spent. The
stopping logic is shown in Figure~22.

\subsubsection{Sub-Group Analyses}

Prespecified subgroup summaries are descriptive and hypothesis-generating only.
Safety and the supportive surgical-quality endpoints are summarized by
resectability category (resectable versus borderline-resectable) and by
daraxonrasib dose level (DL1 through DL3). Given the small sample, subgroup
results are reported as counts and proportions with exact confidence intervals and
are interpreted with explicit caution; no formal test of interaction or of
subgroup difference is performed, and no clinical conclusion is drawn from any
single subgroup.

\subsubsection{Tabulation of Individual Participant Data}

Because the cohort is small and intensively monitored, individual participant data
are tabulated in full. Participant-level listings are provided for demographics
and baseline characteristics; dose level and actual exposure; all adverse events
with onset, severity, seriousness, relationship, and outcome; DLTs and the
resulting escalation decisions; serious AEs and deaths; surgical-quality outcomes;
and the per-procedure device-telemetry and Physical AI event summaries, each keyed
to participant and timepoint. Listings are presented so that every aggregate
statistic can be traced to its contributing records, consistent with the
transparency expected of a first-in-human cyber-physical investigation.

\subsubsection{Exploratory Analyses}

Exploratory analyses are non-confirmatory and are clearly labeled as such.
Progression-free survival (PFS) and overall survival (OS) through the 24-month
follow-up are summarized by the Kaplan-Meier method with the number at risk, and
are interpreted as descriptive given the sample and the absence of a comparator.
The concordance between LLM advisories and the actions ultimately taken by the
operating surgeon, and the sim-to-real gap between the Phase 0 simulation
predictions and the observed intraoperative telemetry, are summarized descriptively
to inform later-phase design and the verification-before-generation calibration.
These exploratory results generate hypotheses for a later-phase study and are not
used to draw efficacy conclusions in this protocol.
